% =============================================================================
% Reporte explicativo (punto de vista de teoria de control)
% HIL de pendulo invertido en TMS320F28379D (LAUNCHXL-F28379D)
%
% COMPILAR:  pdflatex reporte_teoria_control.tex  (dos veces, para indices/refs)
% IMAGENES:  las figuras marcadas "[Insertar imagen]" muestran un recuadro. Para
%            sustituirlas, guarda tu imagen con el nombre indicado en el recuadro
%            (carpeta del .tex o carpeta capturas/) y vuelve a compilar: el
%            documento la usa automaticamente, sin tocar el codigo.
% =============================================================================
\documentclass[11pt,a4paper]{article}

\usepackage[utf8]{inputenc}
\usepackage[T1]{fontenc}
\usepackage[spanish,es-tabla]{babel}
\usepackage{lmodern}
\usepackage{amsmath,amssymb,bm}
\usepackage{booktabs,array,tabularx}
\usepackage{graphicx}
\usepackage{float}
\usepackage{xcolor}
\usepackage{listings}
\usepackage{tikz}
\usetikzlibrary{arrows.meta,positioning,calc}
\usepackage[margin=2.5cm]{geometry}
\usepackage{enumitem}
\usepackage{caption}
\usepackage[hidelinks]{hyperref}

\graphicspath{{./}{capturas/}{../capturas/}{figuras/}}
\setlength{\parskip}{4pt}
\setlength{\emergencystretch}{3em}
\captionsetup{font=small,labelfont=bf}

% ----------------------------- atajos de notacion ----------------------------
\newcommand{\dd}{\mathrm{d}}
\newcommand{\gL}{\tfrac{g}{L}}
\newcommand{\invJ}{\tfrac{1}{J}}
\newcommand{\LSB}{\mathrm{LSB}}

% ------------------- recuadro de imagen pendiente / imagen real --------------
% \espacioimagen{archivo.png}{Pie de figura}{Que debe mostrar la imagen}
\newcommand{\espacioimagen}[3]{%
  \begin{figure}[H]\centering
  \IfFileExists{#1}%
    {\includegraphics[width=.88\linewidth]{#1}}%
    {\fbox{\parbox[c][5.2cm][c]{.84\linewidth}{\centering
       \textbf{[Insertar imagen]}\\[3pt]
       \texttt{\detokenize{#1}}\\[5pt]
       \small #3}}}%
  \caption{#2}
  \end{figure}}

% ----------------------------- estilo de codigo ------------------------------
\definecolor{cgris}{gray}{0.45}
\definecolor{cverde}{rgb}{0.10,0.45,0.15}
\definecolor{cazul}{rgb}{0.10,0.20,0.65}
\lstdefinestyle{codigo}{
  basicstyle=\ttfamily\footnotesize,
  keywordstyle=\color{cazul}\bfseries,
  commentstyle=\color{cverde}\itshape,
  numberstyle=\tiny\color{cgris},
  numbers=left, numbersep=6pt, stepnumber=1,
  breaklines=true, breakatwhitespace=false,
  frame=single, rulecolor=\color{cgris},
  showstringspaces=false, tabsize=2,
  xleftmargin=1.2em, framexleftmargin=1.2em,
  captionpos=b, aboveskip=8pt, belowskip=6pt
}
\lstdefinelanguage{matlabx}{
  morekeywords={for,end,if,else,elseif,while,function,return,true,false,clc,clear,close,all},
  morecomment=[l]{\%}, sensitive=true}
\lstset{style=codigo}

\title{\textbf{Simulaci\'on Hardware-in-the-Loop de un p\'endulo invertido\\
en el TMS320F28379D: fundamento te\'orico de control,\\
implementaci\'on y validaci\'on contra MATLAB}}
\author{Eduardo Pascual P\'erez\\[2pt]
\small Maestr\'ia en Rob\'otica, Universidad Tecnol\'ogica de la Mixteca}
\date{\today}

\begin{document}
\maketitle
\tableofcontents
\newpage

% =============================================================================
\section{Objetivo y alcance}
% =============================================================================
El trabajo implementa un banco de pruebas \emph{Hardware-in-the-Loop} (HIL) en el que
\textbf{el controlador es real} (un PID ejecutado en el DSP TMS320F28379D) y
\textbf{la planta es simulada} (un p\'endulo invertido cuyo estado se integra num\'ericamente
dentro del propio DSP). Planta y controlador se comunican a trav\'es de las
\emph{se\~nales el\'ectricas reales} de los convertidores del chip: la salida del
DAC representa el \'angulo del p\'endulo y la entrada del ADC es lo que ``ve'' el
controlador. Una segunda etapa a\~nade una perturbaci\'on de par introducida mediante
un potenci\'ometro de precisi\'on.

Este documento explica, desde la teor\'ia de control, \emph{qu\'e} se est\'a simulando
(secci\'on~\ref{sec:modelo}), \emph{c\'omo} se controla y qu\'e propiedades tiene el lazo
(secci\'on~\ref{sec:control}), \emph{c\'omo} se discretiz\'o y se implement\'o en el DSP
(secciones~\ref{sec:disc}--\ref{sec:codigo}), \emph{por qu\'e} el an\'alisis en MATLAB se
dise\~n\'o como se dise\~n\'o (secci\'on~\ref{sec:matlab}) y \emph{qu\'e resultados} se
obtuvieron (secci\'on~\ref{sec:resultados}).

\paragraph{Qu\'e se valida con este montaje.} Un HIL no valida el modelo f\'isico (el
p\'endulo es simulado), sino la \textbf{implementaci\'on del controlador}: sus unidades,
su discretizaci\'on, su temporizaci\'on, el efecto de la cuantizaci\'on de 12 bits y de los
convertidores reales. La referencia de comparaci\'on es la misma ley de control
simulada en MATLAB; la diferencia entre ambas son, por construcci\'on, efectos del hardware.

% =============================================================================
\section{Modelo del sistema simulado}\label{sec:modelo}
% =============================================================================
\subsection{P\'endulo invertido con par de entrada}
Se modela una masa puntual $m$ en el extremo de una varilla ligera de longitud $L$,
articulada en la base, con fricci\'on viscosa $b$ y un par de control $u$ aplicado en la
articulaci\'on. El \'angulo $\theta$ se mide respecto a la vertical ($\theta=0$ es el equilibrio
inestable). La perturbaci\'on externa $d$ es un par adicional. La ecuaci\'on de Newton-Euler
para la rotaci\'on es
\begin{equation}
J\,\ddot\theta = m g L \sin\theta - b\,\dot\theta + u + d, \qquad J = mL^2 .
\label{eq:nolineal}
\end{equation}
El t\'ermino $mgL\sin\theta$ es el par gravitatorio y tiene signo \emph{desestabilizante}:
si $\theta>0$ la gravedad aumenta $\theta$. Por eso el sistema en lazo abierto es inestable
y requiere realimentaci\'on.

\subsection{Linealizaci\'on usada en el firmware}
El firmware (y MATLAB) integran la versi\'on linealizada, $\sin\theta\approx\theta$:
\begin{equation}
\ddot\theta = \frac{g}{L}\,\theta - \frac{b}{J}\,\dot\theta + \frac{1}{J}\,(u+d).
\label{eq:lineal}
\end{equation}
En el c\'odigo esto aparece como \texttt{gL*theta - b\_fric*invJ*theta\_dot + invJ*(u + d)}.
La aproximaci\'on es v\'alida mientras $|\theta|$ sea peque\~no: el error relativo de
$\sin\theta\approx\theta$ es $\approx\theta^2/6$, es decir $0.17\,\%$ para la condici\'on inicial
$\theta_0=0.1$~rad y $0.46\,\%$ en el pico m\'as grande observado con perturbaci\'on
($0.167$~rad). Adem\'as, el controlador mantiene $\theta$ cerca de cero, que es justo
donde la linealizaci\'on es mejor. Es una decisi\'on de modelado, no un error del c\'odigo:
cambiarla a $\sin\theta$ ser\'ia inmediato, pero entonces el modelo lineal de la
secci\'on~\ref{sec:control} ya no ser\'ia exacto.

\subsection{Representaci\'on en espacio de estados y polos de lazo abierto}
Con $\bm x=[\theta\;\;\omega]^\top$, $\omega=\dot\theta$:
\begin{equation}
\dot{\bm x}=\underbrace{\begin{bmatrix}0&1\\ g/L & -b/J\end{bmatrix}}_{\bm A}\bm x+
\underbrace{\begin{bmatrix}0\\ 1/J\end{bmatrix}}_{\bm B}(u+d),\qquad y=\theta .
\end{equation}
Los polos son las ra\'ices de $s^2+\tfrac{b}{J}s-\tfrac{g}{L}=0$:
\begin{equation}
s_{1,2}=-\frac{b}{2J}\pm\sqrt{\Big(\frac{b}{2J}\Big)^2+\frac{g}{L}}
\;\Rightarrow\; s_1=+5.447\ \mathrm{s^{-1}},\quad s_2=-6.003\ \mathrm{s^{-1}} .
\end{equation}
Hay un polo real en el semiplano derecho: una perturbaci\'on crece como $e^{5.447\,t}$,
duplic\'andose cada $\ln 2/5.447\approx 0.13$~s. Esto fija los requisitos del lazo: el
controlador debe actuar en una escala de decenas de milisegundos (de ah\'i el muestreo
de la derivada a 100~Hz y el c\'alculo del PID a 100~kHz).

\subsection{Par\'ametros y justificaci\'on}
\begin{table}[H]\centering\small
\begin{tabularx}{\linewidth}{@{}l c l X@{}}
\toprule
Par\'ametro & S\'imbolo & Valor & Papel / por qu\'e importa \\ \midrule
Gravedad & $g$ & 9.81 m/s$^2$ & Fija la inestabilidad: $g/L=32.7\ \mathrm{s^{-2}}$. \\
Longitud & $L$ & 0.3 m & Escala de tiempo: polo inestable $\sqrt{g/L}\approx5.7$~rad/s (constante de tiempo $\approx0.18$~s). \\
Masa & $m$ & 0.1 kg & Junto con $L$ da $J=mL^2=9\times10^{-3}$ kg\,m$^2$ y $mgL=0.294$ N\,m (par gravitatorio por radian). \\
Fricci\'on viscosa & $b$ & 0.005 N\,m\,s/rad & Amortiguamiento natural muy peque\~no: $b/J=0.556\ \mathrm{s^{-1}}$ (casi no ayuda; el amortiguamiento lo aporta $K_D$). \\
Inercia & $J$ & 9$\times10^{-3}$ kg\,m$^2$ & $1/J=111.1$: ganancia de la planta (aceleraci\'on por unidad de par). \\
Paso de integraci\'on & $T_s$ & 10 $\mu$s & Mucho menor que cualquier din\'amica ($\omega T_s\sim10^{-4}$); ver sec.~\ref{sec:disc}. \\
Condici\'on inicial & $\theta_0$ & 0.1 rad & Perturbaci\'on inicial de 5.7$^\circ$: excita el transitorio, dentro de la zona lineal. \\
Rango del ADC/DAC & $\pm\pi/2$ & 0--4095 & Mapea el \'angulo a 12 bits: 1~LSB $=0.767$~mrad. \\
Saturaci\'on del control & $u_{\max}$ & $\pm10$ N\,m & L\'imite del actuador. En las pruebas $|u|_{\max}=0.61$: nunca satur\'o. \\ \bottomrule
\end{tabularx}
\caption{Par\'ametros f\'isicos y num\'ericos del modelo (los mismos en el firmware y en MATLAB).}
\end{table}

\paragraph{Unidades.} Como $J=mL^2$ est\'a en kg\,m$^2$, el par de entrada $u$ (y $d$) est\'a en
N\,m. La condici\'on necesaria de estabilizaci\'on proporcional,
$K_P>mgL=0.294$~N\,m/rad (ec.~\ref{eq:routh}), se cumple con un margen grande
($K_P=5$, es decir 17 veces).

% =============================================================================
\section{Control: PID, estabilidad y rechazo de perturbaciones}\label{sec:control}
% =============================================================================
\subsection{Estructura del controlador}
La referencia es $\theta_{\mathrm{ref}}=0$, de modo que el error es $e=-\theta_{\mathrm{med}}$. La
ley implementada es
\begin{equation}
u = K_P\,e - K_D\,\hat\omega + K_I\!\int e\,\dd t,\qquad
K_P=5,\ K_I=2,\ K_D=0.15 ,
\end{equation}
con tres decisiones de dise\~no que conviene justificar:

\begin{enumerate}[leftmargin=1.6em]
\item \textbf{Derivada sobre la medici\'on} ($-K_D\hat\omega$, no $K_D\dot e$). Con
referencia constante ambas coinciden, pero la forma elegida evita el ``golpe derivativo''
si en el futuro la referencia cambia en escal\'on.
\item \textbf{Derivada lenta y filtrada.} $\hat\omega$ se calcula cada 10~ms como
$\omega_m=(\theta_k-\theta_{k-1})/0.01$ y se filtra con un IIR de primer orden:
\begin{equation}
\hat\omega_k=(1-\alpha)\,\hat\omega_{k-1}+\alpha\,\omega_{m,k},\qquad \alpha=0.386 .
\label{eq:iir}
\end{equation}
La medici\'on llega cuantizada a 0.767~mrad; derivarla a 100~kHz producir\'ia saltos de
$0.767\,\mathrm{mrad}/10\,\mu\mathrm{s}=77$~rad/s. Derivando cada 10~ms, un escal\'on de
1~LSB vale s\'olo $0.0767$~rad/s y, tras el filtro y $K_D$, mueve $u$ en
$0.15\cdot0.386\cdot0.0767\approx4.4\times10^{-3}$. El IIR equivale a una constante de tiempo
$\tau=-T/\ln(1-\alpha)=0.0205$~s (frecuencia de corte $\approx7.8$~Hz).
\item \textbf{Anti-windup por \emph{clamp} del integrador} a $\pm2$ y saturaci\'on de $u$ a
$\pm10$. Evita que el t\'ermino integral crezca sin l\'imite cuando $u$ satura.
\end{enumerate}

\subsection{Funci\'on de transferencia y polinomio caracter\'istico}
Con derivada ideal y la ec.~\eqref{eq:lineal}, sustituyendo $u=-K_P\theta-K_D\dot\theta-K_I\!\int\theta$:
\begin{equation}
\Big(s^2+\tfrac{b+K_D}{J}\,s+\tfrac{K_P}{J}-\tfrac{g}{L}+\tfrac{K_I}{J\,s}\Big)\theta(s)=\tfrac{1}{J}\,d(s)
\;\Rightarrow\;
\frac{\theta(s)}{d(s)}=\frac{\tfrac1J\,s}{s^3+a_2s^2+a_1s+a_0}
\end{equation}
\begin{equation}
a_2=\tfrac{b+K_D}{J}=17.22,\qquad a_1=\tfrac{K_P}{J}-\tfrac gL=522.9,\qquad a_0=\tfrac{K_I}{J}=222.2 .
\end{equation}
\textbf{Criterio de Routh--Hurwitz} (polinomio c\'ubico): estable si $a_2,a_1,a_0>0$ y
$a_2a_1>a_0$. Es decir
\begin{equation}
K_P>mgL,\qquad K_I>0,\qquad \frac{(b+K_D)}{J}\Big(\frac{K_P}{J}-\frac gL\Big)>\frac{K_I}{J}.
\label{eq:routh}
\end{equation}
Num\'ericamente $a_2a_1=9005\gg a_0=222$: estable con holgura. Cada ganancia tiene un
papel claro: $K_P$ debe vencer la gravedad ($a_1$) y fija la rigidez; $K_D$ aporta
amortiguamiento ($a_2$); $K_I$ elimina el error estacionario y determina el polo lento.

\subsection{Polos de lazo cerrado}
\begin{table}[H]\centering\small
\begin{tabular}{@{}lll@{}}\toprule
Modelo & Polos de lazo cerrado & Observaci\'on \\ \midrule
PID con derivada ideal & $-8.40\pm21.10j,\ -0.431$ & $\omega_n=22.7$~rad/s, $\zeta=0.37$ \\
+ filtro de derivada ($\tau=20.5$~ms) & $-8.64\pm26.85j,\ -0.431,\ -31.6$ & $\zeta\approx0.31$ \\ \bottomrule
\end{tabular}
\caption{Polos del lazo cerrado calculados con el modelo lineal (continuo, aproximado).
El muestreo de 100~Hz de la derivada se aproxima como un retardo medio de 5~ms en el an\'alisis
frecuencial de la secci\'on~\ref{sec:margenes}.}
\end{table}
Interpretaci\'on:
\begin{itemize}[leftmargin=1.6em]
\item El \textbf{par complejo} ($\omega_n\approx23$--27~rad/s, $\zeta\approx0.3$--$0.4$) gobierna el
transitorio r\'apido: es subamortiguado, lo que produce el \emph{rebasamiento} que se ve al inicio
de la corrida (el \'angulo cruza el cero y baja a valores negativos antes de volver).
\item El \textbf{polo real lento} $s\approx-a_0/a_1=-0.43$~s$^{-1}$ ($\tau\approx2.3$~s) lo introduce el
integrador. Es el que determina la cola del asentamiento: de ah\'i que, tras una perturbaci\'on, el
\'angulo tarde varios segundos en llegar a $\approx0$ aunque el pico ocurra en 0.12~s.
\end{itemize}

\subsection{Rechazo de perturbaci\'on y error estacionario}
\begin{itemize}[leftmargin=1.6em]
\item \textbf{Con acci\'on integral} $\theta/d$ tiene un cero en $s=0$: ante un escal\'on de $d$,
$\theta(\infty)=0$ y el control compensa exactamente, $u(\infty)=-d$. Es lo que se verifica
experimentalmente con la columna $u+d\approx0$.
\item \textbf{Sin integral} ($K_I=0$) quedar\'ia un error $\theta_\infty=d/(K_P-mgL)=0.2125\,d$:
106~mrad para $d=0.5$. La acci\'on integral es, por tanto, esencial.
\item \textbf{Linealidad.} Sin saturaci\'on, el pico de $\theta$ es proporcional a $d$: la simulaci\'on da
274~mrad por unidad de $d$ (27~mrad para $d=0.1$; 137~mrad para $d=0.5$), siempre en $t\approx0.12$~s.
\item \textbf{Error por offset del ADC.} Si el ADC mide con un sesgo $\delta$ (en LSB), el
controlador lleva a cero la medici\'on, no el \'angulo real, de modo que $\theta_{\mathrm{real}}\to-\delta\cdot0.767$~mrad.
Con la tarjeta a 100~MHz se midi\'o $\delta=-17.6$~LSB. Una captura sin calibrar del repositorio (\texttt{datos\_sin\_perturbacion.csv}) tiene media estacionaria de $13.5$~mrad, que coincide con $17.6\times0.767=13.5$~mrad; es coherente con una corrida a 100~MHz, aunque no consta el reloj con que se tom\'o esa captura. A 200~MHz
se midi\'o $+6.6$~LSB; con la calibraci\'on \texttt{ADC\_CAL\_B}$=6.1$ el residuo es
0.5~LSB ($\approx0.4$~mrad), que coincide con los 0.2--0.5~mrad observados.
\end{itemize}

\subsection{M\'argenes de estabilidad (an\'alisis frecuencial aproximado)}\label{sec:margenes}
Con $L(s)=G(s)\,C(s)$, $G(s)=\dfrac{1/J}{s^2+\frac bJ s-\frac gL}$ y
$C(s)=K_P+\dfrac{K_I}{s}+\dfrac{K_D\,s}{\tau s+1}e^{-sT_d}$ ($\tau=20.5$~ms, $T_d=5$~ms como
retardo equivalente del muestreo a 100~Hz), se obtiene:
\begin{table}[H]\centering\small
\begin{tabular}{@{}ll@{}}\toprule
Cantidad & Valor aproximado \\ \midrule
Frecuencia de cruce $\omega_c$ & 29 rad/s (4.6 Hz) \\
Margen de fase & $22^\circ$ \\
Margen de retardo adicional ($\mathrm{PM}/\omega_c$) & 13 ms \\
Margen de ganancia inferior (en 3.0 rad/s) & $\approx13\times$ \\
Margen de ganancia superior (en 96 rad/s) & $\approx7.2\times$ (17 dB) \\
Pico de sensibilidad $M_s=\max|1/(1+L)|$ & $2.7$ (en 31 rad/s) \\ \bottomrule
\end{tabular}
\caption{M\'argenes del lazo con el modelo aproximado. Son estimaciones: el controlador real es
h\'ibrido (dos velocidades de muestreo) y aqu\'i se aproxima con un filtro y un retardo.}
\end{table}
Un margen de fase de $\approx22^\circ$ y $M_s\approx2.7$ indican un lazo estable pero
\emph{poco amortiguado} ---consistente con $\zeta\approx0.3$ del par dominante---. Esto est\'a
ligado al tema de amortiguamiento que queda por discutir: el principal responsable del
rebasamiento es el retardo y el filtrado de la derivada, no el modelo de la planta.

\espacioimagen{fig_polos_bode.png}{Mapa de polos y diagrama de Bode/Nyquist de $L(s)$.}
{Mapa de polos de lazo abierto y cerrado (ec. de la sec.~3.3) y diagrama de Bode con $\omega_c$, margen de fase y de ganancia marcados. Se puede generar con el script del Ap\'endice~B.}

% =============================================================================
\section{Discretizaci\'on y temporizaci\'on}\label{sec:disc}
% =============================================================================
\subsection{Esquema multi-tasa}
\begin{table}[H]\centering\small
\begin{tabular}{@{}llll@{}}\toprule
Tarea & Periodo & Frecuencia & Contador en el c\'odigo \\ \midrule
ISR: ADC $\rightarrow$ PID $\rightarrow$ planta $\rightarrow$ DAC & 10 $\mu$s & 100 kHz & Timer0 \\
Derivada cruda + filtro IIR & 10 ms & 100 Hz & \texttt{contador\_derivada} (1000) \\
Telemetr\'ia UART & 10 ms & 100 Hz & \texttt{contador\_serial} (1000) \\
Potenci\'ometro (MODO 1) & 10 ms, desfasado 5 ms & 100 Hz & \texttt{contador\_pot} (1000, fase 500) \\ \bottomrule
\end{tabular}
\caption{Tasas de ejecuci\'on. Timer0: 2000 ciclos a 200~MHz $=10\,\mu$s.}
\end{table}

\subsection{Integrador de la planta: Euler--Cromer}
La planta se integra con Euler semi-impl\'icito (Euler--Cromer): primero se actualiza la
velocidad y \emph{con la velocidad nueva} se actualiza la posici\'on,
\begin{equation}
\omega_{k+1}=\omega_k+T_s\Big(\tfrac gL\theta_k-\tfrac bJ\omega_k+\tfrac1J(u_k+d_k)\Big),\qquad
\theta_{k+1}=\theta_k+T_s\,\omega_{k+1}.
\end{equation}
Es el m\'etodo m\'as barato (sin llamadas a funciones, adecuado para una ISR de 10~$\mu$s) y,
a diferencia del Euler expl\'icito, no inyecta energ\'ia artificial en sistemas oscilatorios.
Con $T_s=10\,\mu$s y la din\'amica m\'as r\'apida del lazo en $\approx30$~rad/s, $\omega T_s\approx3\times10^{-4}$:
el error de discretizaci\'on es despreciable frente a la resoluci\'on del ADC.

\subsection{Retardos del lazo HIL}
La salida del DAC se escribe al final de la ISR y el ADC la convierte en la ISR siguiente: hay
un retardo de un periodo ($10\,\mu$s). A la frecuencia de cruce eso son
$29\cdot10^{-5}\approx0.017^\circ$ de fase: despreciable. En cambio, el retardo de la derivada a
100~Hz (efecto de retenedor de orden cero, $\approx5$~ms) s\'i es relevante: consume del orden de
$29\cdot0.005\cdot57.3\approx8^\circ$ del margen de fase y por eso se incluy\'o en el an\'alisis de la
secci\'on~\ref{sec:margenes}.

\subsection{Cuantizaci\'on}
La conversi\'on \'angulo$\leftrightarrow$c\'odigo es lineal:
\begin{equation}
\mathrm{DAC}=\mathrm{round}\!\Big(\tfrac{\theta-\theta_{\min}}{\theta_{\max}-\theta_{\min}}\,4095\Big),\qquad
\theta_{\mathrm{med}}=\theta_{\min}+\frac{\mathrm{ADC}-B}{M}\cdot\frac{\pi}{4095},
\end{equation}
con $\theta_{\min}=-\pi/2$, $B=$\texttt{ADC\_CAL\_B}$=6.1$ LSB y $M=$\texttt{ADC\_CAL\_M}$=1$.
Un LSB equivale a $\pi/4095=0.7672$~mrad. El DAC se \emph{redondea} (no se trunca): truncar
introduc\'ia un sesgo de $-0.5$~LSB ($-0.38$~mrad). La telemetr\'ia se redondea a 1~mrad, as\'i que
el error entre DSP y simulaci\'on tiene un piso del orden de 1~mrad que no es un fallo.

% =============================================================================
\section{Arquitectura del banco HIL}\label{sec:hil}
% =============================================================================
\begin{figure}[H]\centering
\begin{tikzpicture}[
  blk/.style={draw,rectangle,rounded corners=2pt,minimum height=9mm,align=center,font=\small},
  sum/.style={draw,circle,inner sep=1pt,font=\small},
  >={Stealth[length=2.2mm]}]
\node[sum] (s1) at (0,0) {$+$};
\node[blk,minimum width=20mm] (pid) at (3,0) {PID};
\node[sum] (s2) at (5.6,0) {$+$};
\node[blk,minimum width=42mm,font=\scriptsize] (pl) at (9.4,0) {Planta simulada (en la ISR)\\[1pt]$\ddot\theta=\frac gL\theta-\frac bJ\dot\theta+\frac{u+d}{J}$};
\node[blk,minimum width=26mm,font=\scriptsize] (pert) at (5.6,1.7) {$d$: potenci\'ometro\\(pin 29) o escal\'on};
\node[blk,minimum width=20mm,font=\scriptsize] (dac) at (9.4,-2.6) {DAC-B\\(pin 70)};
\node[blk,minimum width=16mm,font=\scriptsize] (cab) at (6.9,-2.6) {Jumper};
\node[blk,minimum width=20mm,font=\scriptsize] (adc) at (4.4,-2.6) {ADC-A0\\(pin 30)};
\node[blk,minimum width=30mm,font=\scriptsize] (der) at (2.2,-1.3) {Derivada 100 Hz\\+ IIR $\alpha{=}0.386$};
\node (ref) at (-1.9,0) {$\theta_{\mathrm{ref}}=0$};
\coordinate (q) at (2.2,-2.6);
\draw[->] (ref) -- (s1);
\draw[->] (s1) -- node[above,font=\scriptsize]{$e$} (pid);
\draw[->] (pid) -- node[above,font=\scriptsize]{$u$} (s2);
\draw[->] (s2) -- (pl);
\draw[->] (pert) -- node[right,font=\scriptsize]{$d$} (s2);
\draw[->] (pl.south -| dac.north) -- node[right,font=\scriptsize]{$\theta$} (dac.north);
\draw[->] (dac) -- (cab);
\draw[->] (cab) -- (adc);
\draw (adc.west) -- node[below,font=\scriptsize]{$\theta_{\mathrm{med}}$} (q);
\draw[->] (q) -- (der.south);
\draw[->] (der.north) -- ++(0,0.45) -| ([xshift=3mm]pid.south);
\draw[->] (q) -- (0,-2.6) -- node[left,font=\scriptsize]{$-$} (s1.south);
\end{tikzpicture}
\caption{Diagrama de bloques del lazo HIL. La rama de la derivada aporta $-K_D\hat\omega$ al PID. Salvo el trayecto f\'isico DAC$\rightarrow$jumper$\rightarrow$ADC, todo el lazo (PID, derivada, planta) se ejecuta como software en la misma ISR del CPU1.}
\end{figure}

\paragraph{Cadena de datos.} (1) La ISR de 100~kHz lee el ADC, calcula el PID e integra la planta;
(2) escribe el \'angulo resultante al DAC-B; (3) cada 10~ms la ISR activa una bandera y el bucle
principal env\'ia por UART (SCI-A, 115200~baud) una l\'inea de texto con enteros separados por comas;
(4) un script de Python lee el puerto, guarda un CSV y mueve el p\'endulo en CoppeliaSim
(visualizaci\'on); (5) los CSV se comparan despu\'es con MATLAB.

\paragraph{Por qu\'e un lazo f\'isico DAC$\rightarrow$ADC.} Cerrar el lazo por los convertidores reales
permite observar, en un banco sencillo, efectos que una simulaci\'on pura no tiene: cuantizaci\'on, offset
y ganancia de los conversores, ruido, y la dependencia del reloj. Los diagn\'osticos de esta tesis (offset
distinto a 100 y 200~MHz, calibraci\'on) surgen exactamente de ello.

\espacioimagen{foto_montaje.png}{Montaje experimental.}
{Fotograf\'ia del LaunchPad con el jumper DAC-B (pin 70) $\rightarrow$ ADC-A0 (pin 30), el potenci\'ometro en el pin 29 y el cable USB a la PC.}

\espacioimagen{coppelia_pendulo.png}{Visualizaci\'on en CoppeliaSim.}
{Captura de la escena de CoppeliaSim con el p\'endulo y la gr\'afica de $\theta$ (y de $d$ en la etapa con perturbaci\'on).}

% =============================================================================
\section{Implementaci\'on en el DSP: c\'odigo y configuraci\'on}\label{sec:codigo}
% =============================================================================
Los fragmentos siguientes son extractos literales de \texttt{control\_hil\_v3.c} y
\texttt{control\_perturbacion.c} (comentarios sin acentos por compatibilidad con el compilador).

\subsection{Reloj del sistema y su efecto}
Con el s\'imbolo predefinido \texttt{\_LAUNCHXL\_F28379D} a nivel de \emph{proyecto}, el SDK configura
SYSCLK$=200$~MHz y LSPCLK$=50$~MHz. Sin \'el, el PLL queda a 100/25~MHz: el Timer0 corre a la
mitad de velocidad (la ISR ya no dura 10~$\mu$s), el divisor de la UART cambia y, adem\'as, cambia el
offset DAC$\rightarrow$ADC (sec.~\ref{sec:resultados}). El c\'odigo lo previene as\'i:
\begin{lstlisting}[language=C,caption={Selecci\'on autom\'atica de LSPCLK y c\'alculo del divisor de baudios.}]
#ifdef _LAUNCHXL_F28379D
  #define LSPCLK_HZ  50000000UL
#else
  #warning "Falta _LAUNCHXL_F28379D: el reloj podria quedar en 100 MHz"
  #define LSPCLK_HZ  25000000UL
#endif
#define BAUD     115200UL
#define BRR_VAL  ((LSPCLK_HZ + 4UL*BAUD) / (8UL*BAUD) - 1UL) // 53 (50 MHz) o 26 (25 MHz)
\end{lstlisting}
El valor $\mathrm{BRR}=\mathrm{round}\big(\mathrm{LSPCLK}/(8\,\mathrm{BAUD})\big)-1$ da $53$ a 50~MHz:
baudios reales $=50\,\mathrm{MHz}/(8\cdot54)=115\,741$ ($+0.47\,\%$, dentro de tolerancia). El valor
anterior, 27 a 25~MHz, daba $-3.1\,\%$ y era el l\'imite de lo tolerable.

\subsection{Par\'ametros del controlador y de la planta}
\begin{lstlisting}[language=C,caption={Constantes del lazo (id\'enticas en MATLAB).}]
#define THETA_INIT  0.1f          // condicion inicial (rad)
#define Ts          0.00001f      // 10 us: paso de la planta y del PID
#define KP          5.0f
#define KI          2.0f
#define KD          0.15f
#define U_MIN       -10.0f
#define U_MAX       10.0f
#define THETA_MIN   -1.570796f    // -pi/2 <-> codigo 0
#define THETA_MAX    1.570796f    // +pi/2 <-> codigo 4095
#define THETA_PER_LSB ((THETA_MAX-THETA_MIN)/4095.0f)   // 0.7672 mrad
#define g 9.81f
#define L 0.3f
#define m 0.1f
#define b_fric 0.005f
#define J    (m*L*L)
#define invJ (1.0f/J)
#define gL   (g/L)
\end{lstlisting}
Se usan \texttt{float} de 32 bits (el F28379D tiene unidad de punto flotante de precisi\'on simple),
con macros en lugar de variables para que el compilador pliegue las constantes y la ISR sea corta.

\subsection{La interrupci\'on principal (n\'ucleo del HIL)}
\begin{lstlisting}[language=C,caption={ISR del Timer0, cada 10 us: ADC $\rightarrow$ PID $\rightarrow$ planta $\rightarrow$ DAC.}]
__interrupt void cpu_timer0_isr(void)
{
    // 1) Leer ADC: convierte lo que el DAC dejo en el pin
    AdcaRegs.ADCINTFLGCLR.bit.ADCINT1 = 1;
    AdcaRegs.ADCSOCFRC1.bit.SOC0 = 1;                 // disparo por software
    volatile int timeout = 1000;
    while ((AdcaRegs.ADCINTFLG.bit.ADCINT1 == 0) && (timeout > 0)) timeout--;
    if (timeout == 0) adc_timeouts++;                 // debe ser 0
    AdcaRegs.ADCINTFLGCLR.bit.ADCINT1 = 1;

    // 2) Codigo -> angulo medido, con calibracion
    Uint16 adc_raw = AdcaResultRegs.ADCRESULT0;
    theta_measured = THETA_MIN + (((float)adc_raw - ADC_CAL_B)*ADC_CAL_INV_M)*THETA_PER_LSB;

    // 3) PID
    float error = 0.0f - theta_measured;
    error_integral += error*Ts;
    if (error_integral >  2.0f) error_integral =  2.0f;     // anti-windup
    if (error_integral < -2.0f) error_integral = -2.0f;

    contador_derivada++;
    if (contador_derivada >= 1000) {                        // cada 10 ms
        theta_dot_meas  = (theta_measured - theta_prev_slow)/0.01f;
        theta_prev_slow = theta_measured;
        contador_derivada = 0;
        theta_dot_filt  = (1.0f-0.386f)*theta_dot_filt + 0.386f*theta_dot_meas;
    }
    u = (KP*error) - (KD*theta_dot_filt) + (KI*error_integral);
    if (u > U_MAX) u = U_MAX;
    if (u < U_MIN) u = U_MIN;

    // 4) Planta (Euler-Cromer)
    float theta_dot_new = theta_dot + Ts*(gL*theta - (b_fric*invJ*theta_dot) + invJ*u);
    float theta_new     = theta + Ts*theta_dot_new;
    if (theta_new > THETA_MAX) theta_new = THETA_MAX;
    if (theta_new < THETA_MIN) theta_new = THETA_MIN;
    theta = theta_new;  theta_dot = theta_dot_new;

    // 5) Angulo -> codigo DAC (lazo de vuelta), con redondeo
    float theta_frac = (theta - THETA_MIN)/(THETA_MAX-THETA_MIN);
    if (theta_frac < 0.0f) theta_frac = 0.0f;
    if (theta_frac > 1.0f) theta_frac = 1.0f;
    dac_code = (Uint16)(theta_frac*4095.0f + 0.5f);
    DacbRegs.DACVALS.all = dac_code;
    ...
    PieCtrlRegs.PIEACK.all = PIEACK_GROUP1;
}
\end{lstlisting}
\paragraph{Lectura paso a paso.}
\begin{enumerate}[leftmargin=1.6em]
\item \emph{Lectura del ADC por software.} \texttt{ADCSOCFRC1} fuerza el inicio de conversi\'on y se
espera la bandera \texttt{ADCINT1}. El \emph{timeout} protege contra un bloqueo y \texttt{adc\_timeouts}
cuenta los fallos: en las pruebas es $0$.
\item \emph{Conversi\'on a \'angulo.} Aqu\'i act\'ua la calibraci\'on $B$/$M$: se resta el offset medido del lazo
DAC$\rightarrow$ADC antes de escalar.
\item \emph{PID.} El orden importa para la comparaci\'on con MATLAB: integral $\rightarrow$ derivada (cada
1000 ISR) $\rightarrow$ $u$ con saturaci\'on. La derivada filtrada $\hat\omega$ s\'olo cambia cada 10~ms.
\item \emph{Planta.} Se usa $u$ reci\'en calculado (sin retardo de c\'omputo: la planta es simulada).
\item \emph{DAC.} La salida se limita a $[0,1]$ antes de convertir a c\'odigo para no desbordar el
\texttt{Uint16}.
\end{enumerate}

\subsection{Configuraci\'on de ADC, DAC y Timer}
\begin{lstlisting}[language=C,caption={Inicializaci\'on de perif\'ericos (extracto).}]
// ADC-A: canal A0 (pin 30) = theta
CpuSysRegs.PCLKCR13.bit.ADC_A = 1;
AdcaRegs.ADCCTL2.bit.PRESCALE = 6;
AdcSetMode(ADC_ADCA, ADC_RESOLUTION_12BIT, ADC_SIGNALMODE_SINGLE); // carga trims
AdcaRegs.ADCSOC0CTL.bit.CHSEL   = 0;     // ADCINA0
AdcaRegs.ADCSOC0CTL.bit.ACQPS   = 14;    // ventana de adquisicion
AdcaRegs.ADCSOC0CTL.bit.TRIGSEL = 0;     // solo por software
AdcaRegs.ADCINTSEL1N2.bit.INT1SEL = 0;   // ADCINT1 = fin de SOC0
AdcaRegs.ADCINTSEL1N2.bit.INT1E   = 1;
// DAC-B (pin 70) con referencia VREFHI
DacbRegs.DACCTL.bit.DACREFSEL = 1;
DacbRegs.DACOUTEN.bit.DACOUTEN = 1;
// Timer0: ISR cada 10 us (200 MHz -> 2000 ciclos)
ConfigCpuTimer(&CpuTimer0, 200, 10);
\end{lstlisting}
\texttt{AdcSetMode()} carga los \emph{trims} de offset del ADC grabados de f\'abrica (necesario para
12 bits seg\'un la errata del dispositivo); sin \'el, el offset es mayor y menos repetible. El valor de
\texttt{ACQPS} no se ha caracterizado experimentalmente: para el canal del potenci\'ometro depende de su
impedancia de fuente y queda como punto a verificar.

\subsection{Arranque limpio y trigger}
La corrida empieza cuando Python env\'ia el car\'acter \texttt{'S'}. La funci\'on \texttt{ResetEstado()} pone
todo el estado en su condici\'on inicial \emph{y} deja el DAC en el c\'odigo de $\theta_0$, de modo que la
primera lectura del ADC ya vea $0.1$~rad:
\begin{lstlisting}[language=C,caption={Arranque coherente con la condici\'on inicial.}]
theta = THETA_INIT;  theta_dot = 0.0f;  theta_prev_slow = THETA_INIT;
theta_dot_filt = 0.0f;  error_integral = 0.0f;  contador_derivada = 0;
dac_code = (Uint16)(((THETA_INIT-THETA_MIN)/THETA_RANGE)*4095.0f + 0.5f);
DacbRegs.DACVALS.all = dac_code;
DELAY_US(200);                       // deja asentar la salida del DAC
\end{lstlisting}
Esta correcci\'on elimin\'o un salto falso de la derivada ($+2.96$~rad/s en la primera trama del firmware
anterior, cuando el DAC arrancaba en 2048, es decir en $\theta=0$).

\subsection{Telemetr\'ia sin \texttt{sprintf}}
El 'int' del C28x es de 16 bits y \texttt{sprintf} provoc\'o un flujo infinito de caracteres '9' que
imped\'ia terminar la primera trama. Se sustituy\'o por un formateador propio que usa solo restas de
potencias de diez:
\begin{lstlisting}[language=C,caption={Formateador entero$\rightarrow$texto y conversi\'on a mrad saturada.}]
static long to_milli(float x){                  // mrad redondeado, saturado a +-999999
    float y = x*1000.0f;
    if (!(y == y)) return 0L;                   // NaN
    if (y >  999999.0f) y =  999999.0f;
    if (y < -999999.0f) y = -999999.0f;
    return (long)(y + ((y >= 0.0f) ? 0.5f : -0.5f));
}
static int fmt_long(char *p, long v){
    static const unsigned long pot[7] = {1000000UL,100000UL,10000UL,1000UL,100UL,10UL,1UL};
    unsigned long a; int n=0,k,started=0;
    if (v < 0) { p[n++]='-'; a=(unsigned long)(-v); } else a=(unsigned long)v;
    for (k=0;k<7;k++){ int d=0; while(a>=pot[k]){a-=pot[k];d++;}
        if (d!=0 || started || k==6){ p[n++]=(char)('0'+d); started=1; } }
    return n;
}
\end{lstlisting}
La l\'inea resultante (\texttt{pos\_mrad,vel\_mrad/s,u\_milli[,pert\_milli]}) cabe en $\approx$25--37 caracteres:
a 115\,741~baud ($\approx86\,\mu$s por car\'acter de 10 bits) son menos de 3.2~ms, bien por debajo de los 10~ms
entre tramas.

\espacioimagen{watch_variables.png}{Vista \emph{Watch} de Code Composer Studio durante una corrida.}
{Captura con \texttt{sysclk\_mhz\_reg}=200, \texttt{isr\_overruns}=0, \texttt{adc\_timeouts}=0, \texttt{start\_flag}=1, \texttt{pot\_raw}, etc.}

\subsection{Etapa con perturbaci\'on}
La perturbaci\'on $d$ entra a la planta sumada al control: \texttt{invJ*(u + pert\_filt)}. Hay tres modos:
0 sin perturbaci\'on, 1 potenci\'ometro y 2 escal\'on programado (repetible, para comparar con MATLAB).

\paragraph{Lectura del potenci\'ometro sin penalizar el lazo.} El canal del potenci\'ometro (ADCINA2, pin~29)
usa un segundo SOC. Se lee a 100~Hz (no a 100~kHz), \emph{desfasado} 5~ms de la derivada y la telemetr\'ia para no
alargar la ISR m\'as pesada. SOC0 ($\theta$) tiene prioridad alta y la ISR espera \emph{ambas} banderas antes de leer
resultados: sin ello, con prioridad por turnos podr\'ia convertirse primero SOC1 y retrasar la lectura de $\theta$.
\begin{lstlisting}[language=C,caption={Conversi\'on del potenci\'ometro a perturbaci\'on (zona muerta, escala sim\'etrica, filtro).}]
static float pot_a_pert(Uint16 raw){
    float x = (float)((int)raw - POT_CENTRO);
    if      (x >  (float)POT_ZONA_MUERTA) x -= (float)POT_ZONA_MUERTA;  // continuidad
    else if (x < -(float)POT_ZONA_MUERTA) x += (float)POT_ZONA_MUERTA;
    else                                  x  = 0.0f;                    // zona muerta
    x = x / (float)(POT_RECORRIDO - POT_ZONA_MUERTA);                   // escala simetrica
    if (x >  1.0f) x =  1.0f;
    if (x < -1.0f) x = -1.0f;
    return x * PERT_MAX;                                                // +-0.5
}
...
if (leer_pot) {
    pot_raw  = AdcaResultRegs.ADCRESULT1;
    pert_cmd = pot_a_pert(pot_raw);
    pert_filt = (1.0f - POT_ALFA)*pert_filt + POT_ALFA*pert_cmd;   // IIR a 100 Hz
}
\end{lstlisting}
\begin{itemize}[leftmargin=1.6em]
\item \emph{Zona muerta de 30 c\'odigos con continuidad:} el potenci\'ometro nunca queda exactamente en el
centro; la zona muerta garantiza $d=0$ con el mando ``en reposo'' y, al restar el ancho de la zona fuera de ella,
$d$ no da un salto al salir.
\item \emph{Escala sim\'etrica} (mismo recorrido a ambos lados) para que $+d$ y $-d$ sean comparables.
\item \emph{Filtro IIR} con $\alpha=0.386$ a 100~Hz ($f_c\approx$10~Hz): suaviza el ruido del potenci\'ometro y los
saltos del mando sin atenuar las perturbaciones relevantes (la din\'amica del lazo est\'a en $\sim$5~Hz).
\item \emph{Calibraci\'on:} el acumulador que estima el offset ($\theta$ medio, \texttt{b\_sugerido}) s\'olo es
v\'alido con $d=0$ y se reinicia cuando hay perturbaci\'on activa.
\end{itemize}

% =============================================================================
\section{An\'alisis en MATLAB: dise\~no y justificaci\'on}\label{sec:matlab}
% =============================================================================
\subsection{Idea general}
MATLAB reproduce \emph{el mismo algoritmo discreto} que corre en el DSP, con el mismo $T_s$, el mismo orden de
operaciones y las mismas constantes, y lo alimenta con los mismos datos de entrada. Despu\'es compara
$\theta$, $\hat\omega$ y $u$ contra la telemetr\'ia del DSP.

\paragraph{Por qu\'e simular el algoritmo discreto y no resolver la ecuaci\'on continua (\texttt{ode45}).}
El objetivo es validar \emph{la implementaci\'on}. Si se compara contra la soluci\'on continua, cualquier
discrepancia mezclar\'ia (i) error de discretizaci\'on, (ii) diferencias de implementaci\'on y (iii) efectos del hardware,
sin poder separarlas. Con el mismo esquema discreto, (i) y (ii) quedan eliminados y la diferencia restante
se atribuye al hardware: cuantizaci\'on del ADC/DAC, redondeo de la telemetr\'ia y temporizaci\'on.

\subsection{Decisiones de dise\~no del script y su justificaci\'on}
Cada decisi\'on responde a una fuente concreta de falsa discrepancia que se encontr\'o en versiones previas:
\begin{enumerate}[leftmargin=1.6em]
\item \textbf{Condiciones iniciales del firmware, no de la primera fila del CSV.}
$\theta=0.1$, $\omega=0$, $\theta_{\mathrm{prev}}=0.1$, $\hat\omega=0$, integral $=0$. La columna ``Velocidad'' del
CSV es $\hat\omega$ (derivada \emph{filtrada de la medici\'on}), no la velocidad real de la planta: usarla como
$\omega_0$ inyectaba $+2.96$~rad/s falsos.
\item \textbf{``Velocidad'' simulada $=\hat\omega$.} El DSP registra $\hat\omega$; comparar contra $\omega$ real
ser\'ia comparar magnitudes distintas (con distinto retardo y filtrado).
\item \textbf{Alineaci\'on temporal.} La fila $k$ del CSV es el estado tras $k\cdot10$~ms. Se simula desde $t=0$
con 1000 pasos de 10~$\mu$s por fila; si se arrancara en la fila 1, todo quedar\'ia corrido 10~ms y el transitorio
(que cambia r\'apido) producir\'ia errores grandes artificiales.
\item \textbf{Redondeo a mrad} (\texttt{round(x*1000)/1000}), igual que \texttt{to\_milli} del firmware, para
comparar peras con peras: el DSP s\'olo transmite 1~mrad de resoluci\'on.
\item \textbf{M\'etricas separadas por r\'egimen.} Transitorio ($t\le0.6$~s) y r\'egimen ($t\ge2$~s) tienen causas de error
distintas (condiciones iniciales vs.\ cuantizaci\'on); un \'unico RMS global las mezclar\'ia y ocultar\'ia cu\'al falla.
\item \textbf{RMS en lugar de error m\'aximo:} el RMS resume todo el registro y es poco sensible a una sola muestra
cuantizada; el pico se reporta aparte.
\end{enumerate}

\begin{lstlisting}[language=matlabx,caption={Lazo de simulaci\'on en MATLAB (misma estructura que la ISR).}]
N_SUB = 1000;                          % 1000 ISR por fila del CSV
for k = 1:N
  for sub = 1:N_SUB
    theta_measured = theta;            % lazo calibrado: sin offset
    error = 0 - theta_measured;
    error_integral = max(min(error_integral + error*Ts, 2), -2);

    contador_derivada = contador_derivada + 1;
    if contador_derivada >= 1000
      theta_dot_meas  = (theta_measured - theta_prev_slow)/0.01;
      theta_prev_slow = theta_measured;
      contador_derivada = 0;
      theta_dot_filt  = 0.614*theta_dot_filt + 0.386*theta_dot_meas;
    end
    u = KP*error - KD*theta_dot_filt + KI*error_integral;
    u = max(min(u, U_MAX), U_MIN);

    theta_dot_new = theta_dot + Ts*(gL*theta - b_fric*invJ*theta_dot + invJ*(u + d));
    theta_new = max(min(theta + Ts*theta_dot_new, THETA_MAX), THETA_MIN);
    theta = theta_new;  theta_dot = theta_dot_new;
  end
  pos_sim(k) = theta;  vel_sim(k) = theta_dot_filt;  u_sim(k) = u;
end
\end{lstlisting}

\subsection{Inyecci\'on de la perturbaci\'on medida}
Con el potenci\'ometro, $d(t)$ no es repetible. En vez de intentar reproducirla, el DSP \emph{registra} la
perturbaci\'on que realmente entr\'o a la planta (cuarta columna, ya filtrada) y MATLAB la usa como entrada
\emph{instante a instante}. As\'i ambos lazos reciben exactamente la misma excitaci\'on, y la comparaci\'on es
v\'alida aunque la perturbaci\'on sea arbitraria. Dos detalles lo hacen exacto:
\begin{itemize}[leftmargin=1.6em]
\item \textbf{Fase de 5~ms.} El DSP actualiza $d$ a la mitad de cada intervalo de telemetr\'ia. La fila $k$ trae el $d$
que entr\'o en la segunda mitad del intervalo $k$; en la primera mitad estaba vigente $d(k-1)$. El script lo replica
(\texttt{usar\_fase\_5ms}).
\item \textbf{Resoluci\'on.} $d$ se transmite redondeada a $10^{-3}$, de ah\'i un error de entrada de hasta
$0.5\times10^{-3}$.
\end{itemize}
\begin{lstlisting}[language=matlabx,caption={Perturbaci\'on vigente en cada paso.}]
if usar_fase_5ms
    if k > 1 && sub <= 500, d = d_dsp(k-1); else, d = d_dsp(k); end
else
    d = d_dsp(k);
end
\end{lstlisting}

\subsection{Verificaci\'on del rechazo: tramos de $d$ constante}
La teor\'ia predice $u+d\to0$ y $\theta\to0$ en cada tramo con $d$ constante. El script detecta tramos de
$\ge3$~s con $|d|$ casi constante y promedia los \'ultimos 1.5~s (cuando el transitorio ya pas\'o):
\begin{lstlisting}[language=matlabx,caption={Detecci\'on de tramos con perturbaci\'on constante.}]
while j < N && abs(d_dsp(j+1) - d_dsp(i)) < 0.01, j = j + 1; end
if t_dsp(j) - t_dsp(i) >= 3.0
    w = max(i, j-150):j;                       % ultimos 1.5 s del tramo
    fprintf(..., mean(d_dsp(w)), mean(pos_dsp(w))*1e3, mean(u_dsp(w)), ...
            mean(u_dsp(w)) + mean(d_dsp(w)));
end
\end{lstlisting}
Se mira $u+d$ y no $u$ solo porque $u+d$ es el par neto que ve la planta: debe ser $\approx0$ (cancelaci\'on
exacta). Es la prueba directa de que el integrador est\'a funcionando.

\subsection{Prueba de la hip\'otesis del arranque}
En las capturas con perturbaci\'on, la primera velocidad filtrada del DSP sale $\approx-2.0$~rad/s cuando el
modelo da $\approx-0.1$~rad/s. Para probar si ese valor explica el transitorio, la opci\'on
\texttt{derivada\_inicial\_csv} reemplaza $\theta_{\mathrm{prev}}$ por el valor que reproduce \emph{esa} primera velocidad:
\begin{lstlisting}[language=matlabx,caption={Replicar el arranque observado de la derivada.}]
if derivada_inicial_csv && primera_derivada
    theta_prev_slow = theta_measured - vel_dsp(1)*0.01/0.386;
end
\end{lstlisting}
Despejando de la ec.~\eqref{eq:iir} con $\hat\omega_0=0$: $\hat\omega_1=\alpha\,(\theta_1-\theta_{\mathrm{prev}})/0.01$,
de donde sale la expresi\'on del c\'odigo. Es una prueba de \emph{consistencia} (si reproduce el transitorio, la
hip\'otesis es plausible); no demuestra por s\'i sola la causa en el hardware, que se confirma leyendo
\texttt{adc\_first[]} y \texttt{theta\_prev\_inicial} en \emph{Watch}.

\espacioimagen{matlab_sin_perturbacion.png}{Comparaci\'on DSP vs MATLAB sin perturbaci\'on.}
{Figura de tres paneles ($\theta$, $\hat\omega$, $u$) de \texttt{comparacion\_sin\_perturbacion.m}, DSP en azul y MATLAB en rojo discontinuo, m\'as el zoom 0--1~s.}

\espacioimagen{matlab_con_perturbacion.png}{Comparaci\'on DSP vs MATLAB con perturbaci\'on.}
{Figura de cuatro paneles ($\theta$, $\hat\omega$, $u$, $d$) de \texttt{comparacion\_con\_perturbacion.m}.}

% =============================================================================
\section{Resultados y discusi\'on}\label{sec:resultados}
% =============================================================================
\subsection{Hallazgos de diagn\'ostico del banco}
\begin{itemize}[leftmargin=1.6em]
\item \textbf{Reloj.} Sin el s\'imbolo \texttt{\_LAUNCHXL\_F28379D} a nivel de proyecto el DSP corr\'ia a
100~MHz (no 200): el Timer0 y la UART quedaban a otra escala. Se verific\'o leyendo
\texttt{sysclk\_mhz\_reg}$=200$ en \emph{Watch}.
\item \textbf{Offset DAC$\rightarrow$ADC.} Cambia con la configuraci\'on de reloj/ADC: $-17.6$~LSB a 100~MHz y
$+6.6$~LSB a 200~MHz. La causa del cambio no se ha verificado. La calibraci\'on por software
(\texttt{ADC\_CAL\_B}$=6.1$) deja un error estacionario de $0.2$--$0.5$~mrad.
\item \textbf{Telemetr\'ia.} El flujo infinito de '9' se eliminó sustituyendo \texttt{sprintf} por \texttt{fmt\_long}.
\item \textbf{Arranque de la derivada.} El firmware anterior arrancaba con el DAC en 2048 ($\theta=0$) y la primera
velocidad sal\'ia $+2.96$~rad/s; se corrigi\'o con \texttt{ResetEstado()}. En el firmware con perturbaci\'on persiste un
artefacto menor (primera velocidad $\approx-2.0$~rad/s) cuya causa en la tarjeta sigue por confirmar.
\end{itemize}

\subsection{Sin perturbaci\'on}
Con el firmware calibrado, en r\'egimen ($t\ge2$~s) la media de $\theta$ medida en una corrida de 20~s fue
$+0.65$~mrad con desviaci\'on est\'andar $0.48$~mrad (capturada con una versi\'on anterior del arranque).
Ambos valores son del orden de 1 LSB ($0.767$~mrad) y de la resoluci\'on de la telemetr\'ia (1~mrad): el lazo
sostiene el p\'endulo en la vertical con un error comparable al del convertidor. La simulaci\'on ideal predice
rebasamiento de $\approx-32$~mrad en $t=0.12$~s y asentamiento a $\pm2$~mrad en $\approx0.57$~s.

\subsection{Con perturbaci\'on}
Corrida de 90~s con el potenci\'ometro moviendo $d$ entre $-0.5$ y $+0.5$. Resultados de la comparaci\'on:
\begin{table}[H]\centering\small
\begin{tabular}{@{}lccc@{}}\toprule
Segmento & RMS $\theta$ (mrad) & RMS $\hat\omega$ (rad/s) & RMS $u$ \\ \midrule
Total & 1.27 & 0.0383 & 0.0080 \\
Transitorio ($t\le0.6$ s) & 8.03 & 0.3263 & 0.0659 \\
Resto ($t\ge2$ s) & 1.09 & 0.0277 & 0.0060 \\ \bottomrule
\end{tabular}
\caption{Error RMS (DSP $-$ MATLAB) con la perturbaci\'on medida como entrada.}
\end{table}
\begin{itemize}[leftmargin=1.6em]
\item \textbf{Picos:} DSP $+167/-95$~mrad contra MATLAB $+169/-93$~mrad (diferencias de 2~mrad $\approx$ 2.6~LSB).
\item \textbf{Error en r\'egimen:} 1.09~mrad RMS, del orden de la resoluci\'on de la telemetr\'ia. El modelo
lineal reproduce la respuesta del DSP a una perturbaci\'on arbitraria.
\item \textbf{Transitorio inicial} (8.03~mrad): es la \'unica zona donde persiste una diferencia. La hip\'otesis del
arranque la explica en simulaci\'on (el RMS baja de 8.07 a $\approx1.0$~mrad al reproducir la primera velocidad
observada), pero \textbf{falta confirmarla en la tarjeta}.
\end{itemize}

\paragraph{Rechazo de la perturbaci\'on (tramos de $d$ constante).}
\begin{table}[H]\centering\small
\begin{tabular}{@{}rrrrrr@{}}\toprule
$t_{ini}$ (s) & $t_{fin}$ (s) & $d$ & $\theta$ (mrad) & $u$ & $u+d$ \\ \midrule
0.0 & 4.9 & $+0.000$ & $-1.03$ & $+0.000$ & $+0.0004$ \\
8.3 & 18.1 & $+0.500$ & $+1.26$ & $-0.500$ & $-0.0007$ \\
21.6 & 31.8 & $-0.313$ & $-2.21$ & $+0.314$ & $+0.0007$ \\
32.4 & 37.5 & $-0.500$ & $-6.43$ & $+0.502$ & $+0.0025$ \\
38.8 & 49.4 & $+0.500$ & $+1.91$ & $-0.501$ & $-0.0008$ \\
53.0 & 68.4 & $+0.003$ & $-0.95$ & $-0.003$ & $+0.0005$ \\
68.9 & 86.6 & $+0.498$ & $-1.23$ & $-0.498$ & $-0.0002$ \\ \bottomrule
\end{tabular}
\caption{Promedios de los \'ultimos 1.5~s de cada tramo con $d$ constante ($\ge3$~s).}
\end{table}
El par neto $u+d$ queda entre $-0.0008$ y $+0.0025$ en todos los tramos: el integrador cancela la perturbaci\'on con
un error menor a 0.5\,\% de $d_{\max}$, y $\theta$ regresa a $\approx0$ (compare con los 106~mrad que dar\'ia un
controlador sin integral para $d=0.5$). El \'unico tramo con $\theta$ notable ($-6.4$~mrad, $d=-0.5$) dur\'o
5.1~s; la simulaci\'on indica que el asentamiento a $\pm2$~mrad ante $d=0.5$ tarda $\approx9$~s (consecuencia del
polo lento en $-0.43$), de modo que lo m\'as probable es que no hubiera terminado de asentarse. Es una explicaci\'on
plausible, pendiente de comprobar sosteniendo $d=-0.5$ durante 15~s o m\'as.

\paragraph{Saturaci\'on.} Con $|u|_{\max}=0.61$ frente al l\'imite de 10, \texttt{u\_sat\_count} no es relevante:
el anti-windup no se activa en estas pruebas, y la hip\'otesis de linealidad del modelo se sostiene.

\espacioimagen{zoom_transitorio.png}{Zoom al transitorio inicial (0--1 s).}
{Superposici\'on DSP/MATLAB mostrando el rebasamiento (DSP $\approx-25$ mrad, MATLAB $\approx-32$ mrad) y su acuerdo despu\'es de $\approx0.4$~s.}

\espacioimagen{rechazo_perturbacion.png}{Rechazo de perturbaci\'on en tramos constantes.}
{$\theta(t)$, $u(t)$ y $d(t)$ de la corrida de 90~s, con los tramos de $d$ constante resaltados.}

\subsection{Discusi\'on}
\begin{itemize}[leftmargin=1.6em]
\item El \textbf{modelo y el PID quedan validados} en el sentido HIL: el DSP reproduce, con diferencias del orden de
la resoluci\'on de los convertidores, la ley de control simulada, tanto sin perturbaci\'on como con una perturbaci\'on
arbitraria medida.
\item Las propiedades predichas por la teor\'ia aparecen en los datos: cancelaci\'on exacta de $d$ ($u+d\approx0$),
linealidad del pico con $d$, y el error estacionario por offset del ADC ($\theta_{\mathrm{real}}=-\delta\cdot0.767$~mrad).
\item El \textbf{amortiguamiento} del par dominante es bajo ($\zeta\approx0.3$--0.4, margen de fase $\approx22^\circ$).
La decisi\'on de modificarlo ---por ejemplo, aumentar $K_D$ o mejorar la derivada--- queda para una discusi\'on
posterior; este documento s\'olo muestra c\'omo se vincula con el filtrado y el retardo de la derivada.
\end{itemize}

\subsection{Limitaciones y pendientes}
\begin{itemize}[leftmargin=1.6em]
\item Confirmar en la tarjeta la causa del artefacto de arranque (\texttt{adc\_first[]}, \texttt{theta\_prev\_inicial}).
\item Captura con escal\'on programado (\texttt{MODO\_PERT}=2) para comparar con una entrada repetible.
\item Margen de ciclos de la ISR (\texttt{isr\_max\_cycles}, l\'imite 2000): medido entre 1378 y 1578 ciclos a 200~MHz, sin \emph{overruns} (capturas del 1 de noviembre y corridas con autocalibraci\'on).
\item Verificar \texttt{ACQPS} del canal del potenci\'ometro seg\'un su impedancia de fuente.
\item Causa del cambio de offset con el reloj: sin determinar. La autocalibraci\'on lo mide y compensa en cada arranque (18 corridas en 3 arranques, dentro de $\pm0.4$~mrad).
\item El an\'alisis de m\'argenes usa un modelo aproximado del controlador h\'ibrido; no es una identificaci\'on experimental.
\end{itemize}

% =============================================================================
\appendix
\section{Variables de inspecci\'on en \emph{Watch}}
% =============================================================================
\begin{table}[H]\centering\small
\begin{tabularx}{\linewidth}{@{}l X@{}}\toprule
Variable & Qu\'e indica \\ \midrule
\texttt{sysclk\_mhz\_reg} & SYSCLK le\'ido de los registros del PLL (debe ser 200). \\
\texttt{isr\_max\_cycles} & Ciclos m\'aximos de la ISR; l\'imite 2000 (10 $\mu$s a 200 MHz). \\
\texttt{isr\_overruns} & Veces que la ISR dur\'o m\'as de 10 $\mu$s (debe ser 0). \\
\texttt{adc\_timeouts} & Conversiones que no terminaron a tiempo (debe ser 0). \\
\texttt{offset\_lsb} & Media de (ADC $-$ DAC) en LSB: offset del lazo f\'isico. \\
\texttt{theta\_mean\_acc} & Media de $\theta$ real en 10~s (error estacionario). \\
\texttt{b\_sugerido} & Valor sugerido de \texttt{ADC\_CAL\_B}. \\
\texttt{pot\_raw} & C\'odigo crudo del potenci\'ometro (para fijar el centro). \\
\texttt{tx\_frames} & Tramas UART enviadas ($\approx100$/s). \\
\texttt{u\_sat\_count} & Veces que $u$ lleg\'o a saturar. \\
\texttt{adc\_first[8]}, \texttt{theta\_prev\_inicial} & Diagn\'ostico del arranque de la derivada. \\ \bottomrule
\end{tabularx}
\end{table}

\section{Script para polos, m\'argenes y figura de la secci\'on 3}
\begin{lstlisting}[language=matlabx,caption={C\'alculo de polos y m\'argenes (requiere Control System Toolbox para \texttt{margin}).}]
g=9.81; L=0.3; m=0.1; b=0.005; J=m*L^2;
KP=5; KI=2; KD=0.15; tau=0.0205; Td=0.005;
s = tf('s');
G = (1/J)/(s^2 + (b/J)*s - g/L);
C = KP + KI/s + KD*s/(tau*s+1)*exp(-Td*s);
Lazo = G*C;
roots([1, (b+KD)/J, KP/J - g/L, KI/J])      % polos con derivada ideal
[Gm,Pm,Wcg,Wcp] = margin(Lazo)             % margenes (el lazo es inestable en OL)
pole(feedback(Lazo,1))                      % polos de lazo cerrado (con retardo)
figure; bode(Lazo); grid on;
figure; nyquist(Lazo); grid on;
\end{lstlisting}
Nota: \texttt{margin} devuelve el margen de ganancia m\'as cercano; con una planta inestable conviene revisar el
diagrama de Nyquist para interpretar si es un margen superior o inferior.

\end{document}
